\chapter{Lange Produkte und globale Zerlegung}
\label{chap:b37-guide-global}

\section{Idempotente Teilprodukte}

\begin{definition}[Teilfolgenprodukte und Blockprodukte]
Für eine Folge $s_1,\ldots,s_m\in S$ mit $m\geq1$ heißt
$s_{i_1}\cdots s_{i_r}$ mit $r\geq1$ und
$1\leq i_1<\cdots<i_r\leq m$ ein \emph{Teilfolgenprodukt}.
Ein \emph{Blockprodukt} ist ein Produkt $s_i s_{i+1}\cdots s_j$
mit $1\leq i\leq j\leq m$.
Die Reihenfolge bleibt in beiden Fällen erhalten; nur bei einer
Teilfolge dürfen dazwischenliegende Glieder ausgelassen werden.
Es sei $q(S)=|S\setminus E(S)|$ die Anzahl der nichtidempotenten Elemente.
\end{definition}

\begin{theorem}[Idempotente Teilfolgen: Gillam--Hall--Williams]
\label{guide:idempotent-subsequence}
Jede Folge $s_1,\ldots,s_m\in S$ mit $m\geq q(S)+1$ besitzt ein
idempotentes Teilfolgenprodukt. Es gibt also
$1\leq i_1<\cdots<i_r\leq m$ mit $r\geq1$ und
\[
 s_{i_1}\cdots s_{i_r}\in E(S).
\]
Die Schranke $q(S)+1$ ist als universelle, allein von der Anzahl der
nichtidempotenten Elemente abhängige Schranke optimal.
\cite{gillam1972,wang2016}
\end{theorem}


\begin{theorem}[Idempotente Blöcke: Hall--Sapir]
\label{guide:idempotent-block}
Jede Folge $s_1,\ldots,s_m\in S$ mit $m\geq 2^{q(S)}$ besitzt ein
idempotentes Blockprodukt:
\[
 \exists\,1\leq i\leq j\leq m:\qquad
 s_i s_{i+1}\cdots s_j\in E(S).
\]
Auch diese Schranke ist als universelle, allein von $q(S)$ abhängige
Schranke optimal. \cite{hallsapir1996}
\end{theorem}

\section{Faktorisierungswälder}

\begin{definition}[Idempotenter Faktorbaum]
\label{def:b37-guide-factor-tree}
Zu einer Folge $s_1,\ldots,s_m\in S$ mit $m\geq 1$ ist ein
\emph{Faktorbaum} ein endlicher geordneter Wurzelbaum, dessen Blätter
von links nach rechts mit $s_1,\ldots,s_m$ beschriftet sind.
Jeder innere Knoten hat mindestens zwei Kinder. Der Wert eines Knotens
ist das Produkt der Blattbeschriftungen unter diesem Knoten in ihrer
gegebenen Reihenfolge.

Der Faktorbaum heißt \emph{idempotent}, wenn an jedem Knoten mit
mindestens drei Kindern alle Kinder denselben Wert $e\in S$ haben und
$e^2=e$ gilt. Seine Höhe ist die größte Anzahl von Kanten auf einem
Weg von der Wurzel zu einem Blatt; ein einzelnes Blatt hat Höhe $0$.
Eine Wahl eines solchen Baums für jede nichtleere endliche Folge in $S$
heißt \emph{Faktorisierungswald}.
\end{definition}

\begin{theorem}[Faktorisierungswaldsatz von Simon]
\label{thm:b37-guide-factor-forest}
Für jede endliche Halbgruppe $S$ und jede Folge
$s_1,\ldots,s_m\in S$ mit $m\geq1$ gibt es einen idempotenten
Faktorbaum der Höhe höchstens $3|S|$.
Die Schranke ist unabhängig von der Länge $m$ des Produkts.
\cite{Simon1990,Kufleitner2008}
\end{theorem}

\section{Unendliche Folgen und Ramsey-Faktorisierung}

\begin{definition}[Verknüpftes Paar]
Ein Paar $(s,e)\in S\times S$ heißt \emph{verknüpft}, wenn
$e^2=e$ und $se=s$ gelten.
\end{definition}

\begin{theorem}[Ramsey-Faktorisierung über einer endlichen Halbgruppe]
\label{guide:ramsey-factorization}
Zu jeder unendlichen Folge $(a_i)_{i\geq1}$ von Elementen aus $S$
gibt es ein verknüpftes Paar $(s,e)$ und ganze Zahlen
\[
 0=k_0<k_1<k_2<\cdots
\]
mit
\[
 a_1\cdots a_{k_1}=s,\qquad
 a_{k_j+1}\cdots a_{k_{j+1}}=e\quad(j\geq1).
\]
\cite{carton2007}
\end{theorem}

\section{Zerlegung nach Krohn und Rhodes}

\begin{definition}[Teilbarkeit von Halbgruppen]
\label{def:b37-guide-division}
Eine Halbgruppe $A$ \emph{teilt} eine Halbgruppe $B$, geschrieben
$A\prec B$, wenn es eine produktabgeschlossene Teilmenge $C\subseteq B$ und einen
surjektiven Halbgruppenhomomorphismus $C\to A$ gibt.
Für diesen Teilbarkeitsbegriff darf $C$ auch leer sein.
\end{definition}

\begin{definition}[Kranzprodukt]
\label{def:b37-guide-wreath}
Für endliche Halbgruppen $A,B$ bezeichnet $A^{B^1}$ die Menge aller
Abbildungen $B^1\to A$. Das \emph{Kranzprodukt} $A\wr B$ ist die
Halbgruppe auf $A^{B^1}\times B$ mit der Verknüpfung
\[
 (f,b)(g,c)=(h,bc),
 \qquad h(x)=f(x)\,g(xb)\quad(x\in B^1).
\]
Ein \emph{iteriertes Kranzprodukt} entsteht durch endlich viele
Anwendungen dieser Konstruktion mit festgelegter Klammerung.
\end{definition}

\begin{definition}[Bausteine der Primzerlegung]
\label{def:b37-guide-prime-blocks}
Die Halbgruppe $U=\{1,r_0,r_1\}$ ist durch
\[
 1u=u1=u\quad(u\in U),\qquad
 r_i r_j=r_j\quad(i,j\in\{0,1\})
\]
festgelegt. Eine \emph{einfache Gruppe} ist eine nichttriviale Gruppe,
deren einzige Normaluntergruppen die triviale Untergruppe und die
Gruppe selbst sind. Eine Untergruppe $N$ von $G$ heißt dabei normal,
wenn $gNg^{-1}=N$ für alle $g\in G$ gilt.
Eine endliche einfache Gruppe $G$ heißt
\emph{einfacher Gruppenteiler von $S$}, wenn $G\prec S$ gilt.
\end{definition}

\begin{theorem}[Primzerlegungssatz von Krohn--Rhodes]
\label{thm:b37-guide-krohn-rhodes}
Jede endliche Halbgruppe $S$ teilt ein endliches iteriertes Kranzprodukt,
dessen Faktoren Kopien von $U$ und endliche einfache Gruppenteiler
von $S$ sind. \cite{KrohnRhodes1965,RhodesSteinberg2009}
\end{theorem}

\begin{theorem}[Einfache Gruppenteiler von Kranzprodukten]
\label{thm:b37-guide-prime-group-divisors}
Es seien $T_1,\ldots,T_r$ endliche Halbgruppen mit $r\geq1$ und
$W$ ein aus diesen Faktoren gebildetes iteriertes Kranzprodukt
mit beliebiger festgelegter Klammerung.
Für jede endliche einfache Gruppe $G$ gilt
\[
 G\prec W\quad\Longrightarrow\quad
 G\prec T_i\text{ für mindestens ein }i\in\{1,\ldots,r\}.
\]
\cite[Satz~4.1.45]{RhodesSteinberg2009}
\end{theorem}

\clearpage
\section{Rees-Zerlegung aus Gruppenteilern}

\begin{definition}[Monoid und Rees-Erweiterung]
Eine Halbgruppe mit einem neutralen Element heißt \emph{Monoid}.
Für Monoide $N,L$ und eine Abbildung $\rho:N\to L$ ist die
\emph{Rees-Erweiterung} $\operatorname{Rees}(N,L,\rho)$ das Monoid
auf der disjunkten Vereinigung
\[
 N\mathbin{\dot\cup}(N\times L\times N)
\]
mit der Multiplikation
\[
\begin{aligned}
 n\cdot n'&=nn',\\
 n\cdot(a,\ell,b)&=(na,\ell,b),\\
 (a,\ell,b)\cdot n&=(a,\ell,bn),\\
 (a,\ell,b)\cdot(c,k,d)&=(a,\ell\rho(bc)k,d).
\end{aligned}
\]
Hierbei liegen $n,n',a,b,c,d$ in $N$ und $\ell,k$ in $L$.
\end{definition}

\begin{theorem}[Endliche Rees-Zerlegung nach Diekert--Walter]
\label{thm:b37-guide-rees-decomposition}
Für jede endliche Halbgruppe $S$ gibt es mit $M=S^1$ eine Folge
endlicher Monoide $R_1,\ldots,R_r$ mit
\[
 1\leq r\leq 2^{|M|}-1,\qquad M\prec R_r,
\]
so daß für jedes $j\in\{1,\ldots,r\}$ einer der folgenden Fälle gilt:
\begin{enumerate}
 \item $R_j$ ist eine Gruppe, die $M$ teilt.
 \item Es gibt $i,h\in\{1,\ldots,j-1\}$ und eine Abbildung
 $\rho_j:R_i\to R_h$ mit
 \[
 R_j=\operatorname{Rees}(R_i,R_h,\rho_j).
 \]
\end{enumerate}
\cite[Lemma~1 und Proposition~14]{DiekertWalter2016}
\end{theorem}

\clearpage
\section{Komplexität der Zerlegung}

\begin{definition}[Aperiodizität und Krohn--Rhodes-Komplexität]
\label{def:b37-guide-complexity}
Eine endliche Halbgruppe heißt \emph{aperiodisch}, wenn jede ihrer
Untergruppen aus genau einem Element besteht.
Die \emph{Krohn--Rhodes-Komplexität} $c_{\mathrm{KR}}(S)$ ist die kleinste
Anzahl nichttrivialer Gruppenfaktoren in einem endlichen iterierten
Kranzprodukt $W$ mit $S\prec W$, dessen Faktoren endliche Gruppen
oder endliche aperiodische Halbgruppen sind.
Die Gruppenfaktoren müssen dabei nicht einfach sein.
Für $S=\varnothing$ setzen wir $c_{\mathrm{KR}}(S)=0$.
\end{definition}

\begin{definition}[Gruppentiefe]
Die \emph{Gruppentiefe} $\delta(S)$ ist die größte Anzahl von Klassen
in einer strikt aufsteigenden Kette
\[
 J_1<J_2<\cdots<J_r
\]
von $\mathcal J$-Klassen, von denen jede eine Untergruppe mit mindestens
zwei Elementen enthält. Gibt es keine solche Klasse, setzen wir
$\delta(S)=0$. Gezählt werden die Klassen, nicht die Vergleichsschritte.
\end{definition}

\begin{theorem}[Tiefenzerlegungssatz von Rhodes--Tilson]
\label{thm:b37-guide-depth}
Für jede endliche Halbgruppe $S$ gilt
\[
 c_{\mathrm{KR}}(S)\leq\delta(S).
\]
\cite{margolis2025}
\end{theorem}

\begin{theorem}[Berechenbarkeit der Krohn--Rhodes-Komplexität]
\label{thm:b37-guide-complexity-decidability}
Es gibt einen Algorithmus, der aus der Multiplikationstafel jeder
endlichen Halbgruppe $S$ die Zahl $c_{\mathrm{KR}}(S)$ berechnet.
\cite{margolis2024,margolis2025}
\end{theorem}

\section{Komplexität homomorpher Bilder}

\begin{theorem}[Fundamentallemma der Komplexität von Rhodes]
\label{thm:b37-guide-fundamental-complexity}
Für jede endliche Halbgruppe $S$ und jeden surjektiven Homomorphismus
$f:S\to T$, dessen Einschränkung auf jede Untergruppe von $S$
injektiv ist, gilt
\[
 c_{\mathrm{KR}}(S)\leq c_{\mathrm{KR}}(T).
\]
\cite[Satz~4.9.1]{RhodesSteinberg2009}
\end{theorem}
